%% toc: P1 6 6.7
\section{Soundness, completeness, consistency, and their distinct implications}
\label{sec:p01-w02-p01-019-consistency}

The definitions of soundness and completeness in \cref{sec:levels-03} compare
two relations on the same formulas: derivability \(\Gamma\derives_{\mathcal
C}\varphi\) in a specified calculus \(\mathcal C\), and semantic consequence
\(\Gamma\entails\varphi\) in a specified semantics. The direction check
recorded there separates the two implications
\cite[ch.~``Levels, Objects, and Consequence'', definitions ``Soundness'' and
``Completeness'', warning ``Direction check'']{FrozenFormalLogicConsistencyComparison}.
A third notion, consistency, is a property of a premise set relative to a
calculus. The handbook defines it separately for \(\mathsf{ND}_{\mathsf{CPL}}\)
(\cref{sec:prop-metatheory-02}), for \(\mathsf{ND}_{\mathsf{FOL}}\)
(\cref{sec:fol-metatheory-01}) and for theories in arithmetic
(\cref{sec:first-incompleteness-02}). This section gives the general
definition for an arbitrary specified calculus. It then states exactly what each
of soundness and completeness implies about consistency, and shows by
examples that no implication beyond the stated ones holds.

The inputs are a language with formula set \(F\), a semantics
\(\mathcal S\) with permitted data \(\mathcal D_{\mathcal S}\) and
satisfaction relation \(d\entails_{\mathcal S}\varphi\)
(\cref{def:p01-w02-p01-016-datum,def:p01-w02-p01-016-satisfaction}), the
induced consequence relation \(\entails_{\mathcal S}\) and satisfiability
(\cref{def:p01-w02-p01-016-satisfiable}), and a calculus \(\mathcal C\) with
derivability \(\derives_{\mathcal C}\) as defined in \cref{sec:levels-03}. The
conclusions are as follows.
\begin{itemize}
  \item Consistency has three forms, absolute, negation and falsum
        consistency. They coincide under named structural hypotheses
        (\cref{thm:p01-w02-p01-019-forms}), which
        \(\mathsf{ND}_{\mathsf{CPL}}\) satisfies
        (\cref{thm:p01-w02-p01-019-ndcpl}).
  \item Soundness makes every satisfiable set consistent and, given
        reflexivity and the classical negation clause, is equivalent to that
        property (\cref{thm:p01-w02-p01-019-soundness}).
  \item Completeness makes every consistent set satisfiable and, for a
        calculus with reductio under the classical negation clause, is
        equivalent to that property (\cref{thm:p01-w02-p01-019-completeness}).
  \item For an adequate calculus, consistency is satisfiability, and
        deductive completeness of a set is a semantic condition
        (\cref{thm:p01-w02-p01-019-adequate}).
  \item Explicit calculi separate the notions
        (\cref{thm:p01-w02-p01-019-three-calculi,thm:p01-w02-p01-019-separation}).
        Consistency of a calculus does not give soundness, soundness with
        explosion does not give completeness, and soundness and completeness
        of a calculus do not give deductive completeness of a premise set.
\end{itemize}
The results and their proofs are editorial completions built on the
inherited definitions. They are not attributed to an external presentation.

\subsection{Standing data and structural properties}
\label{sec:p01-w02-p01-019-standing}

Throughout, \(F\) is a nonempty set of formulas closed under a unary
connective \(\neg\). The connective may be primitive, as in classical
propositional logic, or defined, as in intuitionistic logic, where
\(\neg A\) abbreviates \(A\to\bottom\)
\cite[ch.~``Intuitionistic Logic'', sec.~``Language, calculus, and Kripke
structures'']{FrozenFormalLogicConsistencyComparison}. The letters
\(\Gamma,\Delta\) range over subsets of \(F\), and
\(\varphi,\psi,\chi\) over members of \(F\). The semantics \(\mathcal S\) and
the calculus \(\mathcal C\) are fixed by the context, and the subscripts are
dropped when no other semantics or calculus is in play.

\begin{hbdefinition}{def:p01-w02-p01-019-structural}{Structural properties of a calculus and of negation}
A calculus \(\mathcal C\) over \(F\) has:
\begin{enumerate}
  \item \emph{reflexivity} iff \(\varphi\in\Gamma\) implies
        \(\Gamma\derives_{\mathcal C}\varphi\);
  \item \emph{explosion} iff, for all \(\Gamma,\varphi,\psi\),
        \(\Gamma\derives_{\mathcal C}\varphi\) and
        \(\Gamma\derives_{\mathcal C}\neg\varphi\) imply
        \(\Gamma\derives_{\mathcal C}\psi\);
  \item \emph{reductio} iff, for all \(\Gamma,\varphi\), if some \(\psi\)
        satisfies \(\Gamma\cup\{\neg\varphi\}\derives_{\mathcal C}\psi\) and
        \(\Gamma\cup\{\neg\varphi\}\derives_{\mathcal C}\neg\psi\), then
        \(\Gamma\derives_{\mathcal C}\varphi\);
  \item when \(\bottom\in F\), the \emph{falsum properties} iff
        (B1) \(\Gamma\derives_{\mathcal C}\bottom\) implies
        \(\Gamma\derives_{\mathcal C}\psi\) for every \(\psi\), and
        (B2) \(\Gamma\derives_{\mathcal C}\varphi\) and
        \(\Gamma\derives_{\mathcal C}\neg\varphi\) imply
        \(\Gamma\derives_{\mathcal C}\bottom\).
\end{enumerate}
A semantics \(\mathcal S\) has \emph{exclusive negation} iff no datum
satisfies both \(\varphi\) and \(\neg\varphi\). It has the \emph{classical
negation clause} iff \(d\entails_{\mathcal S}\neg\varphi\) holds exactly when
\(d\entails_{\mathcal S}\varphi\) fails, for every datum \(d\) and formula
\(\varphi\). The classical clause implies exclusive negation. The three
semantics of \cref{def:p01-w02-p01-016-datum} have the classical clause.
\index{explosion (of a calculus)}\index{reductio (property of a calculus)}
\end{hbdefinition}

Reflexivity is a property of the calculus, not of the definition of
derivability. In \(\mathsf{ND}_{\mathsf{CPL}}\) it is witnessed by the
one-node derivation. It fails in a calculus whose derivations admit no
assumption leaves. Monotonicity
(\(\Gamma\subseteq\Delta\) and \(\Gamma\derives_{\mathcal C}\varphi\) imply
\(\Delta\derives_{\mathcal C}\varphi\)) needs no hypothesis. A derivation
whose undischarged assumptions lie in \(\Gamma\) has them in \(\Delta\)
(clause 0 of \cref{thm:p01-w02-p01-019-forms}).

\subsection{Consistency relative to a calculus}
\label{sec:p01-w02-p01-019-forms}

\begin{hbdefinition}{def:p01-w02-p01-019-consistency}{Consistency, triviality and deductive completeness relative to a calculus}
Let \(\mathcal C\) be a specified calculus over \(F\) and
\(\Gamma\subseteq F\).
\begin{enumerate}
  \item \(\Gamma\) is \emph{\(\mathcal C\)-trivial} iff
        \(\Gamma\derives_{\mathcal C}\psi\) for every \(\psi\in F\), and
        \emph{absolutely \(\mathcal C\)-consistent} otherwise.
  \item \(\Gamma\) is \emph{\(\mathcal C\)-inconsistent} iff there is
        \(\varphi\in F\) with \(\Gamma\derives_{\mathcal C}\varphi\) and
        \(\Gamma\derives_{\mathcal C}\neg\varphi\), and
        \emph{\(\mathcal C\)-consistent} otherwise. This is
        \emph{negation consistency}. Unqualified ``consistent'' in this
        section means it, as in the inherited definition for theories
        \cite[ch.~``Diagonalization and the First Incompleteness Theorem'',
        definition ``Consistency'']{FrozenFormalLogicConsistencyComparison}.
  \item When \(\bottom\in F\), \(\Gamma\) is
        \emph{\(\bottom\)-consistent} iff
        \(\Gamma\notderives_{\mathcal C}\bottom\). This is the form of the
        inherited definitions for \(\mathsf{ND}_{\mathsf{CPL}}\) and
        \(\mathsf{ND}_{\mathsf{FOL}}\) \cite[ch.~``Soundness, Completeness,
        Compactness, and Resolution'', definition ``Consistency''; ch.~``Henkin
        Completeness and Compactness'', definition ``First-order
        consistency'']{FrozenFormalLogicConsistencyComparison}.
  \item The calculus \(\mathcal C\) is \emph{consistent} (absolutely, by
        negation, or with respect to \(\bottom\)) iff the empty set is.
  \item \(\Gamma\) is \emph{deductively complete for \(\mathcal C\)} iff,
        for every \(\varphi\in F\), \(\Gamma\derives_{\mathcal C}\varphi\) or
        \(\Gamma\derives_{\mathcal C}\neg\varphi\).
\end{enumerate}
\index{consistency!relative to a calculus}\index{trivial set of formulas}
\index{deductively complete set}
\end{hbdefinition}

Clause 5 extends to arbitrary formula sets and calculi the property of
theories fixed in \cref{def:theory-completeness}
\cite[ch.~``Quantifier Elimination and Categoricity'', definition
``Deductive completeness, semantic completeness, and model
completeness'']{FrozenFormalLogicConsistencyComparison}. Every notion in the
definition is relative to \(\mathcal C\). The same set can be consistent for
one calculus and inconsistent for another, as \cref{thm:p01-w02-p01-019-three-calculi}
shows for the set \(\{\neg P\}\).

\begin{hblemma}{thm:p01-w02-p01-019-forms}{proved=annex:p01-w02-p01-019-forms}{Comparison of the forms of consistency}
Let \(\mathcal C\) be a calculus over \(F\).
\begin{enumerate}
  \setcounter{enumi}{-1}
  \item If \(\Gamma\subseteq\Delta\) and
        \(\Gamma\derives_{\mathcal C}\varphi\), then
        \(\Delta\derives_{\mathcal C}\varphi\). Hence every subset of a
        consistent set is consistent, and the same holds for absolute and
        \(\bottom\)-consistency.
  \item Every \(\mathcal C\)-consistent set is absolutely
        \(\mathcal C\)-consistent.
  \item If \(\mathcal C\) has explosion, every absolutely
        \(\mathcal C\)-consistent set is \(\mathcal C\)-consistent, so the
        two forms coincide.
  \item If \(\mathcal C\) has reductio, it has explosion.
  \item If \(\bottom\in F\) and \(\mathcal C\) has the falsum properties, then
        \(\mathcal C\) has explosion. A set is then \(\bottom\)-consistent iff
        it is absolutely consistent iff it is consistent.
\end{enumerate}
\end{hblemma}

\begin{hblemma}{thm:p01-w02-p01-019-ndcpl}{proved=annex:p01-w02-p01-019-ndcpl}{The structural properties of \(\mathsf{ND}_{\mathsf{CPL}}\)}
The calculus \(\mathsf{ND}_{\mathsf{CPL}}\) of \cref{sec:prop-deduction-01}
has reflexivity, the falsum properties and reductio. Consequently, for every
\(\Gamma\subseteq\mathsf{Fm}_{\mathsf{CPL}}\), the conditions
\(\Gamma\notderives_{\mathsf{ND}_{\mathsf{CPL}}}\bottom\) (the definition of
\cref{sec:prop-metatheory-02}), absolute
\(\mathsf{ND}_{\mathsf{CPL}}\)-consistency and
\(\mathsf{ND}_{\mathsf{CPL}}\)-consistency are equivalent.
\end{hblemma}

\subsection{What soundness gives and what completeness gives}
\label{sec:p01-w02-p01-019-implications}

\begin{hbtheorem}{thm:p01-w02-p01-019-soundness}{proved=annex:p01-w02-p01-019-soundness}{Soundness and consistency}
Let \(\mathcal C\) be sound for \(\mathcal S\).
\begin{enumerate}
  \item If \(\mathcal S\) has exclusive negation, every
        \(\mathcal S\)-satisfiable set is \(\mathcal C\)-consistent, and so
        absolutely \(\mathcal C\)-consistent. If \(\bottom\in F\) and no
        datum satisfies \(\bottom\), every satisfiable set is also
        \(\bottom\)-consistent.
  \item If \(\mathcal S\) has exclusive negation and
        \(\mathcal D_{\mathcal S}\neq\varnothing\), the calculus
        \(\mathcal C\) is consistent.
  \item Equivalently to clause 1, a derivation of some \(\varphi\) and of
        \(\neg\varphi\) from \(\Gamma\) certifies that \(\Gamma\) is
        unsatisfiable.
\end{enumerate}
Conversely, suppose \(\mathcal S\) has the classical negation clause and
\(\mathcal C\) has reflexivity. If every \(\mathcal S\)-satisfiable set is
\(\mathcal C\)-consistent, then \(\mathcal C\) is sound. Under these two
hypotheses, soundness is therefore equivalent to consistency of every
satisfiable set.
\end{hbtheorem}

\begin{hbtheorem}{thm:p01-w02-p01-019-completeness}{proved=annex:p01-w02-p01-019-completeness}{Completeness and consistency}
\begin{enumerate}
  \item For every semantics \(\mathcal S\), if \(\mathcal C\) is complete for
        \(\mathcal S\), then every absolutely \(\mathcal C\)-consistent set
        is \(\mathcal S\)-satisfiable. In particular every
        \(\mathcal C\)-consistent set is satisfiable.
  \item Suppose \(\mathcal S\) has the classical negation clause and
        \(\mathcal C\) has reductio. If every \(\mathcal C\)-consistent set is
        \(\mathcal S\)-satisfiable, then \(\mathcal C\) is complete for
        \(\mathcal S\).
  \item Hence, for \(\mathcal S\) with the classical negation clause and
        \(\mathcal C\) with reductio, the following are equivalent:
        \(\mathcal C\) is complete; every \(\mathcal C\)-consistent set is
        satisfiable; every absolutely \(\mathcal C\)-consistent set is
        satisfiable.
\end{enumerate}
\end{hbtheorem}

Clause 1 is the model-existence form of completeness. The inherited proof of
\cref{thm:prop-completeness} uses clause 2: it shows that
\(\Gamma\cup\{\neg\varphi\}\) is consistent whenever
\(\Gamma\notderives_{\mathsf{ND}_{\mathsf{CPL}}}\varphi\), and then it builds a
valuation satisfying that consistent set. Reductio cannot be weakened to
explosion in clause 2. The calculus \(\mathcal C_{\mathrm{ex}}\) of
\cref{thm:p01-w02-p01-019-three-calculi} is sound and explosive, and every
consistent set is satisfiable for it, but it is not complete.

\begin{hbcorollary}{thm:p01-w02-p01-019-adequate}{proved=annex:p01-w02-p01-019-adequate}{Consistency and deductive completeness for an adequate calculus}
Let \(\mathcal S\) have the classical negation clause, and let \(\mathcal C\)
be sound and complete for \(\mathcal S\).
\begin{enumerate}
  \item For every set, being \(\mathcal C\)-consistent, being absolutely
        \(\mathcal C\)-consistent and being \(\mathcal S\)-satisfiable are
        equivalent.
  \item \(\mathcal C\) is consistent iff
        \(\mathcal D_{\mathcal S}\neq\varnothing\).
  \item \(\Gamma\) is \(\mathcal C\)-consistent and deductively complete for
        \(\mathcal C\) iff \(\Gamma\) is satisfiable and any two data that
        satisfy \(\Gamma\) satisfy exactly the same formulas.
  \item Let \(\mathcal S\) be the first-order semantics of a signature
        \(\Lang\) and \(T\) a set of \(\Lang\)-sentences. Then \(T\) is
        \(\mathcal C\)-consistent and, for every \(\Lang\)-sentence
        \(\sigma\), \(T\derives_{\mathcal C}\sigma\) or
        \(T\derives_{\mathcal C}\neg\sigma\), iff \(T\) has a model and any
        two models of \(T\) are models of exactly the same sentences.
\end{enumerate}
\end{hbcorollary}

Clause 3 quantifies over all formulas. For first-order data it is too strong
to describe theories, because two assignments in one model can disagree on an
open formula such as \(x=y\). Clause 4 restricts the quantifier to sentences,
which is the restriction made in \cref{def:theory-completeness}. The
corollary needs no structural hypothesis on \(\mathcal C\). Soundness and completeness determine
\(\derives_{\mathcal C}\) as the relation \(\entails_{\mathcal S}\), and all
four clauses are statements about that relation.

\subsection{Separating the notions}
\label{sec:p01-w02-p01-019-separations}

The examples use calculi in which no rule discharges an assumption. For them
derivability has a closure description. That description gives exact
derivability conditions without a search through trees.

\begin{hbdefinition}{def:p01-w02-p01-019-rule-calculus}{Calculus given by rule instances}
Let \(R\) be a set of pairs \((\Pi,\chi)\) with \(\Pi\subseteq F\) finite and
\(\chi\in F\). A \emph{derivation in \(\mathcal C_R\)} is a finite rooted tree
whose nodes are labelled by formulas. Each node is of one of two kinds.
It is an \emph{assumption leaf}, which is never discharged, or it is a
\emph{rule node} for some \((\Pi,\chi)\in R\), labelled \(\chi\), whose
children are labelled by the members of \(\Pi\), one child for each member.
A rule node with \(\Pi=\varnothing\) is a leaf, and it is not an assumption.
\(\Gamma\derives_{\mathcal C_R}\chi\) holds iff some derivation has root
label \(\chi\) and all its assumption leaves labelled by members of
\(\Gamma\), as in \cref{sec:levels-03}. A set \(X\subseteq F\) is
\emph{\(R\)-closed} iff \(\Pi\subseteq X\) and \((\Pi,\chi)\in R\) imply
\(\chi\in X\).
\end{hbdefinition}

\begin{hblemma}{thm:p01-w02-p01-019-least-closed}{proved=annex:p01-w02-p01-019-least-closed}{Derivability as least closed superset}
For every \(R\) and \(\Gamma\), the set
\(\{\chi:\Gamma\derives_{\mathcal C_R}\chi\}\) is \(R\)-closed, contains
\(\Gamma\), and is included in every \(R\)-closed superset of \(\Gamma\). In
particular \(\mathcal C_R\) has reflexivity.
\end{hblemma}

\begin{hbdefinition}{def:p01-w02-p01-019-three-calculi}{Three test calculi on the classical propositional language}
Let \(F=\mathsf{Fm}_{\mathsf{CPL}}\) with the valuation semantics of
\cref{sec:prop-syntax-semantics-02}, and let \(P,Q\) be distinct specified
propositional variables. Define:
\begin{enumerate}
  \item \(\mathcal C_0\), the calculus with \(R=\varnothing\);
  \item \(\mathcal C_{\mathrm{ex}}\), with
        \(R_{\mathrm{ex}}=\{(\Pi,\chi):\Pi\text{ finite and unsatisfiable},\
        \chi\in F\}\);
  \item \(\mathcal C_P\), with
        \(R_P=\{(\Pi,\chi):\Pi\text{ finite and }\Pi\cup\{P\}\entails\chi\}\).
\end{enumerate}
\end{hbdefinition}

Each side condition is decidable. It concerns a finite set of formulas, and
only the finitely many variables occurring in it matter
(\cref{thm:p01-w02-p01-016-finite-check}). So all three are calculi with
decidable rule sets. They are built only to test the implications and are not
proposed as logics.

\begin{hbtheorem}{thm:p01-w02-p01-019-three-calculi}{proved=annex:p01-w02-p01-019-three-calculi}{Derivability and the properties of the three test calculi}
For all \(\Gamma\subseteq\mathsf{Fm}_{\mathsf{CPL}}\) and \(\chi\):
\begin{enumerate}
  \item \(\Gamma\derives_{\mathcal C_0}\chi\) iff \(\chi\in\Gamma\). The
        calculus \(\mathcal C_0\) is sound and has reflexivity. It lacks
        explosion and is not complete. The set \(\{P,\neg P\}\) is absolutely
        \(\mathcal C_0\)-consistent, \(\mathcal C_0\)-inconsistent and
        unsatisfiable.
  \item \(\Gamma\derives_{\mathcal C_{\mathrm{ex}}}\chi\) iff
        \(\chi\in\Gamma\) or \(\Gamma\) is unsatisfiable. The calculus
        \(\mathcal C_{\mathrm{ex}}\) is sound, has reflexivity, the falsum
        properties and explosion, and its consistent sets are exactly the
        satisfiable sets. It lacks reductio and is not complete:
        \(\{P\}\entails P\lor Q\) but
        \(\{P\}\notderives_{\mathcal C_{\mathrm{ex}}}P\lor Q\).
  \item \(\Gamma\derives_{\mathcal C_P}\chi\) iff
        \(\Gamma\cup\{P\}\entails\chi\). The calculus \(\mathcal C_P\) is
        complete, has reflexivity, the falsum properties, explosion and
        reductio, and is consistent. Every \(\mathcal C_P\)-consistent set is
        satisfiable. It is not sound: \(\derives_{\mathcal C_P}P\) but
        \(\not\entails P\). The satisfiable set \(\{\neg P\}\) is
        \(\mathcal C_P\)-inconsistent.
\end{enumerate}
\end{hbtheorem}

\begin{hbcorollary}{thm:p01-w02-p01-019-separation}{proved=annex:p01-w02-p01-019-separation}{The three notions are independent}
Over the classical propositional language and its valuation semantics:
\begin{enumerate}
  \item A consistent calculus need not be sound, even if it is complete and
        has reductio. Consistency of the calculus is the single instance
        \(\Gamma=\varnothing\) of the condition in the converse part of
        \cref{thm:p01-w02-p01-019-soundness}, and it does not imply the
        condition for all satisfiable sets (\(\mathcal C_P\)).
  \item A sound calculus need not be complete, even if it is explosive and
        all its consistent sets are satisfiable (\(\mathcal C_{\mathrm{ex}}\)).
        A complete calculus need not be sound (\(\mathcal C_P\)).
  \item Without explosion, absolute consistency is strictly weaker than
        consistency, and completeness fails while soundness holds
        (\(\mathcal C_0\)).
  \item Soundness and completeness of a calculus do not make a consistent set
        deductively complete. The empty set is
        \(\mathsf{ND}_{\mathsf{CPL}}\)-consistent and neither
        \(\derives_{\mathsf{ND}_{\mathsf{CPL}}}P\) nor
        \(\derives_{\mathsf{ND}_{\mathsf{CPL}}}\neg P\) holds. For each
        valuation \(v\), the set
        \(\Gamma_v=\{p\in\mathsf{Prop}:v\entails p\}\cup
        \{\neg p:p\in\mathsf{Prop},\ v\not\entails p\}\) is
        \(\mathsf{ND}_{\mathsf{CPL}}\)-consistent and deductively complete.
\end{enumerate}
\end{hbcorollary}

The comparison can be summarized as follows. Soundness is a property of
\(\mathcal C\) over all premise sets. It is equivalent, under the classical
negation clause and reflexivity, to consistency of every satisfiable set. The
consistency of \(\mathcal C\) is only the instance of that property at the
empty set. Completeness is equivalent to satisfiability of every consistent
set when \(\mathcal C\) has reductio, and under explosion alone it is
stronger. Consistency and deductive completeness are properties of one
premise set relative to one calculus. The first concerns whether the set
derives too much, and the second whether it settles every formula. They are
answered by the semantics only after soundness and completeness of the
calculus are known (\cref{thm:p01-w02-p01-019-adequate}).

\subsection{A worked comparison in classical propositional natural deduction}
\label{sec:p01-w02-p01-019-worked}

\begin{hbexample}{def:p01-w02-p01-019-worked}{Deciding consistency through the semantics}
Let \(\Gamma=\{P\to Q,\ P,\ \neg Q\}\) and \(\Gamma'=\{P\to Q,\ \neg Q\}\) in
\(\mathsf{Fm}_{\mathsf{CPL}}\).

\emph{\(\Gamma\) is inconsistent, shown syntactically.} From the assumptions
\(P\to Q\) and \(P\), the rule \(\to E\) gives \(Q\). From \(Q\) and the
assumption \(\neg Q\), the rule \(\neg E\) gives \(\bottom\). The resulting
derivation has open assumptions \(P\to Q,P,\neg Q\in\Gamma\). Hence
\(\Gamma\derives_{\mathsf{ND}_{\mathsf{CPL}}}Q\) and
\(\Gamma\derives_{\mathsf{ND}_{\mathsf{CPL}}}\neg Q\) (the latter by
reflexivity), and \(\Gamma\derives_{\mathsf{ND}_{\mathsf{CPL}}}\bottom\). By
clause 3 of \cref{thm:p01-w02-p01-019-soundness}, \(\Gamma\) is
unsatisfiable. A direct check agrees: a valuation satisfying \(P\) and
\(P\to Q\) satisfies \(Q\) and so fails \(\neg Q\).

\emph{\(\Gamma'\) is consistent, shown semantically.} The valuation \(v\) with
\(v(P)=v(Q)=\mathbf F\) satisfies \(P\to Q\) and \(\neg Q\). By clause 1 of
\cref{thm:p01-w02-p01-019-soundness} and
\cref{thm:prop-nd-soundness}, \(\Gamma'\) is
\(\mathsf{ND}_{\mathsf{CPL}}\)-consistent, and by
\cref{thm:p01-w02-p01-019-ndcpl}
\(\Gamma'\notderives_{\mathsf{ND}_{\mathsf{CPL}}}\bottom\). This
nonderivability claim covers every derivation. It is established by one
valuation, not by a search through derivations.

\emph{Completeness supplies a derivation.} Every valuation satisfying
\(\Gamma'\) fails \(Q\) and therefore fails \(P\), so
\(\Gamma'\entails\neg P\). By \cref{thm:prop-completeness} a derivation
exists. One such derivation assumes \(P\) with label \(1\), applies
\(\to E\) with \(P\to Q\) to obtain \(Q\), applies \(\neg E\) with \(\neg Q\)
to obtain \(\bottom\), and closes with \(\neg I\) discharging label \(1\).

\emph{\(\Gamma'\) is not deductively complete.} Let \(S\) be a third
propositional variable. The valuations \(v\) above and \(v'\), which agrees
with \(v\) except \(v'(S)=\mathbf T\), both satisfy \(\Gamma'\) and disagree on
\(S\). By clause 3 of \cref{thm:p01-w02-p01-019-adequate}, \(\Gamma'\) derives
neither \(S\) nor \(\neg S\).
\end{hbexample}

\subsection{Limits of the comparison}
\label{sec:p01-w02-p01-019-limits}

\begin{hbwarning}{def:p01-w02-p01-019-limits}{Consistency is relative, and the classical reduction is a hypothesis}
\begin{enumerate}
  \item ``Consistent'' without a named calculus has no fixed meaning. The set
        \(\{\neg P\}\) is \(\mathsf{ND}_{\mathsf{CPL}}\)-consistent and
        \(\mathcal C_P\)-inconsistent
        (\cref{thm:p01-w02-p01-019-three-calculi}).
  \item Showing that a calculus is consistent, for example by finding one
        nonderivable formula, does not show that it is sound
        (\cref{thm:p01-w02-p01-019-separation}, clause 1).
  \item In a calculus without explosion, a set that derives a formula and its
        negation can still fail to derive some other formula, as
        \(\{P,\neg P\}\) does in \(\mathcal C_0\). The forms of consistency
        must then be named separately.
  \item The converse halves of
        \cref{thm:p01-w02-p01-019-soundness,thm:p01-w02-p01-019-completeness}
        use the classical negation clause through clause 1 of
        \cref{thm:p01-w02-p01-016-unsat-reduction}. For the intuitionistic
        Kripke semantics that reduction fails, and every calculus that is
        sound and complete for that semantics lacks reductio
        (\cref{thm:p01-w02-p01-019-intuitionistic}). The forward halves,
        satisfiable implies consistent under soundness and consistent implies
        satisfiable under completeness, remain valid, because intuitionistic
        forcing has exclusive negation.
  \item Deductive completeness of a set is not completeness of a calculus.
        \(\mathsf{ND}_{\mathsf{CPL}}\) is complete, yet the consistent empty
        set is not deductively complete for it
        (\cref{thm:p01-w02-p01-019-separation}, clause 4).
\end{enumerate}
\end{hbwarning}

\begin{hbtheorem}{thm:p01-w02-p01-019-intuitionistic}{proved=annex:p01-w02-p01-019-intuitionistic}{Failure of the classical reduction for intuitionistic forcing}
Take as data the pairs \((\mathcal M,w)\) of an intuitionistic Kripke
structure and a world, with forcing \(\mathcal M,w\forces A\) as
satisfaction, and with \(\neg A\) defined as \(A\to\bottom\)
(\cref{sec:intuitionistic-logic-01}). Then
\(\entails_{\mathsf{IPL}}\) is the consequence relation of
\cref{sec:levels-03} for these data. Let \(p\) be a propositional variable.
\begin{enumerate}
  \item Negation is exclusive: no world forces both \(A\) and \(\neg A\).
  \item No world of any structure forces \(\neg(p\lor\neg p)\), so
        \(\{\neg(p\lor\neg p)\}\) is unsatisfiable.
  \item \(\varnothing\not\entails_{\mathsf{IPL}}p\lor\neg p\).
  \item If \(\mathcal C\) is sound and complete for this semantics, then
        \(\{\neg(p\lor\neg p)\}\) is \(\mathcal C\)-inconsistent while
        \(\varnothing\notderives_{\mathcal C}p\lor\neg p\). So \(\mathcal C\)
        lacks reductio.
\end{enumerate}
\end{hbtheorem}

The inherited chapter states soundness and strong completeness of
\(\derives_{\mathsf{IPL}}\) for this semantics and supplies a proof sketch
\cite[ch.~``Intuitionistic Logic'', theorem ``Soundness and strong
completeness of \(\mathsf{IPL}\)'']{FrozenFormalLogicConsistencyComparison}.
Clause 4 applies to \(\derives_{\mathsf{IPL}}\) on the strength of that
theorem. Clauses 1--3 and the conditional form of clause 4 are proved in full
in the annex.

\subsection{Relationships and sources}
\label{sec:p01-w02-p01-019-sources}

The calculus-specific definitions of consistency in
\cref{sec:prop-metatheory-02} and \cref{sec:fol-metatheory-01} are the
\(\bottom\)-form of \cref{def:p01-w02-p01-019-consistency}. For
\(\mathsf{ND}_{\mathsf{CPL}}\) they coincide with the other two forms by
\cref{thm:p01-w02-p01-019-ndcpl}. The definition for theories in
\cref{sec:first-incompleteness-02} is the negation form, and its
equivalence with \(T\notderives\bottom\) there is an instance of clause 4 of
\cref{thm:p01-w02-p01-019-forms} for a calculus with the falsum properties. Omega-consistency
and \(\Sigma^0_1\)-soundness, defined in the same section, are properties of
theories relative to the structure \(\mathcal N\), not soundness of a calculus.
\Cref{def:theory-completeness} contrasts deductive and semantic completeness
of a theory and identifies them through soundness and completeness of the
first-order calculus. \Cref{thm:p01-w02-p01-019-adequate} is the general form
of that identification. The proof of \cref{thm:prop-completeness} follows the
route of clause 2 of \cref{thm:p01-w02-p01-019-completeness}.

The soundness and completeness definitions and the direction check are those
of the frozen inherited manuscript, reproduced on the active route in
\cref{sec:levels-03}
\cite[ch.~``Levels, Objects, and Consequence'', definitions ``Soundness'' and
``Completeness'', warning ``Direction check'']{FrozenFormalLogicConsistencyComparison}.
Satisfiability and the reduction of consequence to unsatisfiability are taken
from \cref{sec:p01-w02-p01-016-models}. The intuitionistic semantics is the
one of the inherited chapter \cite[ch.~``Intuitionistic Logic'', definitions
``Intuitionistic Kripke frame and structure'' and ``Intuitionistic semantic
consequence'']{FrozenFormalLogicConsistencyComparison}. The general
definition of consistency, the structural properties, the test calculi, all
results of this section and their proofs are editorial.
